\chapter*{You Can Just Do Things}
\addcontentsline{toc}{chapter}{You Can Just Do Things}

In the source code of pandas---the Python library used by millions of data scientists---there's a comment in \texttt{groupby.pyx}:

\begin{lstlisting}[language=python]
# Running stats update based on RunningStats::Push from
# https://www.johndcook.com/blog/skewness_kurtosis/
\end{lstlisting}

In OpenTSDB, the time-series database that powers monitoring at scale, there's a Java comment:

\begin{lstlisting}[language=java]
/**
\begin{itemize}
\item This implementation is based upon a paper by John D. Cook,
\item which itself is based upon a method that goes back to a 1962
\item paper by B. P. Welford
\end{itemize}
 */
\end{lstlisting}

In the official TensorFlow documentation, there's a note in the autodiff guide:

\begin{blogquote}
Note: Don't actually invert the matrix.
\end{blogquote}

These citations: buried in production code. Embedded in official documentation---trace back to a WordPress blog that has been publishing since January 2008. One mathematician, writing about what interested him, day after day, for seventeen years.

This book is about what that looks like.

\bigskip

Simple numbers: 5,233 posts. 6,533 days. An average of one post every 1.25 days. No breaks longer than eleven days. No sabbaticals, no pivots to video, no platform migrations. Just writing.

But numbers don't capture influence. Cook's gamma function implementation has been ported to ten programming languages. His explanation of "don't invert that matrix" became the canonical reference that TensorFlow links to. His running statistics algorithm powers skewness and kurtosis calculations in pandas groupby operations---meaning data scientists computing higher moments on grouped DataFrames are running code that traces back to his blog.

Infrastructure, not fame or brand.

\bigskip

Why does this matter?

Because it demonstrates something rare: the ability to just do things. Write about what interests you. Publish without permission. Build utility over time. Become genuinely useful without virality, platform, or strategy.

Cook's blog appears in university syllabi at MIT, Stanford, and Carnegie Mellon. When a programmer searches "gamma function python" at 2 AM, trying to debug something urgent, they find his explanation. When a graduate student needs an intuitive explanation of Bayesian inference, they find his post. When an engineer needs to understand why naive variance computation is numerically unstable, they find his running statistics tutorial. And then that algorithm ends up in their production code.

Some reach can be measured: 5,881 citations on Google Scholar. Stack Overflow answers that have reached 5.7 million users. Over 1,200 submissions to Hacker News, accumulating 25,000 upvotes. Seventeen specialized Twitter accounts with 500,000 combined followers. A single interpolation calculator page---one of more than eighty standalone tools---has generated over 600,000 search clicks in sixteen months.

Real influence, though, can't be counted. It lives in code comments, in understanding that went uncited, in problems solved at 2 AM that never got recorded.

\bigskip

Not a biography. Cook's personal life isn't the subject.

A data portrait of sustained intellectual work---an attempt to see the shape of that influence through counting, measuring, and visualizing. How many posts about prime numbers? When did cryptography become a major topic? How has the average post length changed?

Data reveals patterns. Patterns suggest a story. An arc emerges from scattered consultant's notebook to mathematical institution: from short posts about PowerShell bugs in 2008 to deep explorations of elliptic curve cryptography in 2025.

It asks what happens when someone just keeps going.

\bigskip

When the blog launched in January 2008, people had blogrolls. Google Reader existed. Twitter asked "what are you doing?" instead of optimizing for engagement.

While the internet reinvented itself three times over, johndcook.com kept doing one thing: math, computing, statistics, the occasional digression into music theory or typography. No ads. No viral moments. No controversy.

Ecology has a term for this: climax species---the organism that remains after initial colonizers have come and gone, after boom-and-bust cycles have played out, after the ecosystem has stabilized.

Viral sensations aren't the internet's climax species. Quiet sites that just keep being useful are.

Consider the comparison. Joel Spolsky's \textit{Joel on Software} shaped a generation of developers, then he co-founded Stack Overflow and largely stopped blogging. Jeff Atwood's \textit{Coding Horror} ran for years before the same pivot. Steve Yegge wrote legendary posts at Google, then went quiet. Paul Graham's essays remain influential but arrive sporadically. These are the famous names in technical blogging---and none of them kept going the way Cook has. They built platforms, started companies, moved on. Cook just kept writing. Seventeen years, same site, same format, same focus on being useful. The rarity isn't the quality. It's the duration.

\bigskip

Seventeen years in three acts.

First, the beginning: 2008 to 2012, when the blog was a consultant's notebook, scattered, practical, exploratory. Second, the crucible: 2013 to 2017, when practice was tested and voice clarified. Third, the institution: 2018 to 2025, when the blog became infrastructure, cited in documentation, ported across languages, running in production systems worldwide.

Data tells this story. Data also misses things. A final chapter acknowledges what counting can't capture: the reader who finally understood, the code that was never attributed, the influence that left no trace.

A map. The territory is larger.

\bigskip

\textit{All statistics in this book reflect the blog as of November 27, 2025. The archive continues to grow.}

